% ==============================================================================
% Fichier     : chapitres/00_Conclusion.tex
% Titre       : Conclusion Générale, Cours de Hacking Éthique
% Auteur      : MINKA MI NGUIDJOÏ Thierry Emmanuel

% ==============================================================================

\chapter*{Conclusion Générale}
\addcontentsline{toc}{chapter}{Conclusion Générale}
\markboth{Conclusion Générale}{Conclusion Générale}

\begin{flushright}
    \itshape
    <<~La sécurité n'est pas un produit, c'est un processus.~>>\\[0.5ex]
    \small--- Bruce Schneier
\end{flushright}

\vspace{1.5ex}

% ==============================================================================
\section*{Ce que Vous Avez Appris}
% ==============================================================================

Au terme de ce parcours offensif, vous avez traversé l'intégralité de la
chaîne d'attaque qui structure le métier de pentester, de la toute première
requête \sigle{DNS} passive jusqu'à la remise du rapport contradictoire avec
plan de remédiation priorisé. Ce cheminement révèle une vérité fondamentale
que seule la pratique peut enseigner : \textit{une infrastructure n'est
jamais plus solide que son maillon le plus faible}, et ce maillon n'est
presque jamais là où on l'attend.

Vous avez compris que le hacking éthique n'est pas une collection de
\og{}tricks\fg{} techniques, mais une \textbf{discipline intellectuelle
structurée} : observer avant d'agir, documenter chaque pas, mesurer l'impact
réel avant de l'exposer, et toujours replacer la technique dans son contexte
métier. La valeur d'un pentest se mesure moins au nombre de shells obtenus
qu'à la clarté des recommandations transmises à la direction.

% ==============================================================================
\section*{Synthèse du Parcours d'Attaque LiZbiZ}
% ==============================================================================

\begin{boxresume}
\textbf{Le scénario LiZbiZ a illustré, de bout en bout, un cycle d'attaque
complet en 8 phases, avec mapping MITRE ATT\&CK :}

\begin{enumerate}
    \item \textbf{Reconnaissance \sigle{OSINT}} \hfill \textit{TA0043}\\
          Découverte passive : 47 emails internes via Google Dork
          (\texttt{filetype:xls "employés"}), compte \texttt{admin-it}
          via métadonnées ExifTool, Apache~2.2.22 via Shodan, sous-domaines
          internes via transfert \sigle{AXFR} et Certificate Transparency logs.

    \item \textbf{Scanning actif et énumération} \hfill \textit{TA0043}\\
          Cartographie réseau complète : détection Apache~2.2.22 non patché,
          Windows Server 2012~R2 exposé port~445, \sigle{LDAP} anonymous bind,
          3 comptes de service avec \sigle{SPN} (cibles Kerberoasting),
          politique sans verrouillage de compte.

    \item \textbf{Phishing opérationnel} \hfill \textit{T1566.002}\\
          Campagne GoPhish sur 47 cibles, taux de clic 34\%, 11 credentials
          capturés dont \texttt{admin-it@lizbiz.com / LiZbiZ2024!}.

    \item \textbf{Exploitation web} \hfill \textit{T1190, T1059.007}\\
          Injection SQL (\texttt{' OR '1'='1' --}) sur le formulaire de
          connexion, extraction de 47 comptes. \sigle{XSS} stored sur
          l'intranet, vol du cookie de session administrateur. \sigle{LFI}
          + Log Poisoning → shell \texttt{www-data} sur le serveur \sigle{DMZ}.

    \item \textbf{Exploitation système} \hfill \textit{T1210}\\
          EternalBlue (\sigle{MS17-010}) sur Windows Server 2012~R2 —
          shell \texttt{NT AUTHORITY\textbackslash{}SYSTEM} obtenu en 28~secondes.
          Payload \texttt{windows/x64/meterpreter/reverse\_tcp}.

    \item \textbf{Post-exploitation et mouvement latéral} \hfill \textit{T1548, T1003, T1550}\\
          Escalade \texttt{sudo vim} → \texttt{root} (serveur Linux).
          Tunnel Chisel \sigle{DMZ} → \sigle{LAN}. Mimikatz
          \texttt{sekurlsa::logonpasswords} → hash + mot de passe en clair
          \texttt{admin-it}. Pass-the-Hash vers le Contrôleur de Domaine.

    \item \textbf{Compromission Active Directory} \hfill \textit{T1558.003, T1003.006, T1558.001}\\
          Kerberoasting → mot de passe \texttt{svc\_backup} cracké (hashcat).
          DCSync → dump complet du domaine incluant hash \texttt{krbtgt}.
          Golden Ticket forgé → persistance illimitée sur \texttt{lizbiz.local}.

    \item \textbf{Rapport et remédiation} \hfill \textit{N/A}\\
          20 fiches de vulnérabilités avec scores \sigle{CVSS v3.1} et
          références \sigle{MITRE ATT\&CK}. Plan de remédiation priorisé.
          Débriefing en 3 temps (exécutif, technique, plan d'action).
\end{enumerate}
\end{boxresume}

% ==============================================================================
\section*{Compétences Consolidées}
% ==============================================================================

Ce parcours a forgé des compétences directement transférables en milieu
professionnel :

\begin{itemize}
    \item \textbf{Rigueur méthodologique}, L'application des standards \sigle{PTES},
          \sigle{NIST SP 800-115} et \sigle{MITRE ATT\&CK} distingue le
          professionnel du \textit{script-kiddie}. Chaque technique est
          documentée, horodatée et rattachée à une tactique connue des
          équipes défensives.

    \item \textbf{Conscience juridique}, Savoir précisément où s'arrête
          le mandat et où commence le délit, dans les deux juridictions
          étudiées (France et Cameroun). La chaîne de custody des preuves
          est une compétence autant juridique que technique.

    \item \textbf{Communication biface}, Rédiger un Executive Summary
          compréhensible par un \sigle{PDG} et une fiche de vulnérabilité
          exploitable par un développeur : les deux formats sont aussi
          importants que les exploits eux-mêmes.

    \item \textbf{Réflexivité défensive}, Comprendre l'attaque pour
          concevoir la contre-mesure. Pour chaque technique offensive
          présentée, une contre-mesure défensive a été identifiée :
          Kerberoasting → mots de passe longs sur les comptes de service ;
          Golden Ticket → rotation du compte \texttt{krbtgt} ;
          fileless → Credential Guard, surveillance comportementale.

    \item \textbf{Vision hybride}, L'architecture LiZbiZ (on-premise +
          cloud) reflète la réalité des \sigle{SI} modernes. Le pentester
          contemporain doit maîtriser autant les attaques \sigle{AD}
          classiques que les \sigle{SSRF} vers les endpoints \sigle{IMDS}
          cloud (169.254.169.254) et les misconfigurations \sigle{IAM}.
\end{itemize}

% ==============================================================================
\section*{Les Vecteurs Émergents : Ce que le Cours Ne Couvre Pas Encore}
% ==============================================================================

La cybersécurité offensive est un domaine en évolution permanente. Plusieurs
vecteurs, mentionnés en modules optionnels dans ce cours, méritent une
attention particulière dans la pratique professionnelle de 2025--2026 :

\begin{itemize}
    \item \textbf{Cloud Security Posture}, Les misconfigurations \sigle{AWS}/Azure/\sigle{GCP}
          (buckets S3 publics, politiques \sigle{IAM} excessives, \sigle{SSRF}
          vers l'\sigle{IMDS}) représentent une part croissante des
          compromissions réelles. Le passage de l'on-premise au cloud
          ne réduit pas la surface d'attaque, il la déplace.

    \item \textbf{Supply Chain Attacks}, SolarWinds, Log4j, MoveIT ont
          démontré qu'un seul composant tiers compromis peut propager une
          attaque à des milliers d'organisations simultanément. L'audit
          des dépendances logicielles (\sigle{SBOM}) devient une compétence
          indispensable.

    \item \textbf{Adversarial \sigle{AI}}, L'utilisation de modèles de
          langage pour automatiser la génération de prétextes de phishing
          hyper-personnalisés, l'obfuscation de code et la découverte de
          vulnérabilités logicielles redéfinit les capacités offensives à
          la portée d'acteurs peu qualifiés.

    \item \textbf{Zero Trust Architecture}, Face à l'évolution des
          menaces, le modèle périmétrique traditionnel cède la place au
          \textit{Zero Trust} : \og{}ne jamais faire confiance, toujours
          vérifier\fg{}. Comprendre ce modèle est indispensable pour
          concevoir des tests d'intrusion pertinents dans ces
          environnements.
\end{itemize}

% ==============================================================================
\section*{La Complémentarité Audit / Pentest}
% ==============================================================================

Ce cours clôt une séquence pédagogique commencée avec le Cours d'Audit des
Systèmes d'Information. La complémentarité entre ces deux disciplines est
profonde et fondamentale :

\begin{center}
\begin{tabular}{p{4cm}p{4.5cm}p{4.5cm}}
\toprule
\textbf{Dimension} & \textbf{Audit \sigle{SI}} & \textbf{Pentest} \\
\midrule
Objet & Conception des contrôles & Résistance des contrôles \\
Méthode & Revue documentaire & Simulation d'attaque \\
Posture & Évaluation & Confrontation \\
Livrable & Rapport de conformité & Rapport de compromission \\
Question & \og{}Les contrôles existent-ils ?\fg{} & \og{}Les contrôles résistent-ils ?\fg{} \\
\bottomrule
\end{tabular}
\end{center}

\medskip
L'un sans l'autre est incomplet. Un audit sans pentest évalue des contrôles
dont l'effectivité n'est pas prouvée. Un pentest sans audit préalable manque
le contexte organisationnel qui donne sens aux vulnérabilités découvertes.
C'est précisément cette dualité, évaluer et tester, qui structure la
pratique du \sigle{RSSI} contemporain.

\begin{boxinfo}
\textbf{Certifications recommandées pour aller plus loin :}
\begin{itemize}
    \item \textbf{\sigle{OSCP}} (\textit{Offensive Security Certified
          Professional}, Offensive Security), la référence mondiale du
          pentest technique : examen pratique de 24h, accès non guidé à
          plusieurs machines.
    \item \textbf{\sigle{PNPT}} (\textit{Practical Network Penetration
          Tester}, TCM Security), orienté pentest complet et reporting,
          accessible aux débutants, reconnu en entreprise.
    \item \textbf{\sigle{CEH}} (\textit{Certified Ethical Hacker},
          EC-Council), largement reconnu en entreprise, particulièrement
          dans les appels d'offres publics africains.
    \item \textbf{\sigle{GPEN}} (GIAC Penetration Tester, SANS) —
          orienté méthodologie et reporting, reconnu dans les contrats
          gouvernementaux.
    \item \textbf{\sigle{CRTO}} (\textit{Certified Red Team Operator},
          Zero Point Security), spécialisé Active Directory, infrastructure
          C2 (Cobalt Strike, Havoc) et \sigle{OPSEC} avancée.
    \item \textbf{\sigle{CISA} / \sigle{CISM}} (ISACA), pour ancrer
          la compétence offensive dans la gouvernance \sigle{SI} et la
          gestion des risques ; particulièrement valorisés dans les
          organisations publiques et les institutions financières.
\end{itemize}
\end{boxinfo}

% ==============================================================================
\section*{Mot Final}
% ==============================================================================

Attaquer pour défendre est une posture intellectuellement exigeante. Elle
requiert de comprendre les systèmes en profondeur, d'anticiper les erreurs
humaines, de modéliser les comportements adverses, et de tout cela,
d'extraire des recommandations utiles, honnêtes et actionnables.

Ce que ce cours vous a, en dernier lieu, appris à faire, c'est \textbf{lire
une infrastructure comme un attaquant et l'expliquer comme un défenseur}.
Cette double lecture est rare. Elle est précieuse. Elle est ce que les
organisations recherchent dans un \sigle{RSSI} de qualité.

Le pentester professionnel est, en définitive, un \textbf{avocat du diable
au service de la sécurité} : il incarne temporairement l'adversaire pour
mieux armer l'organisation contre lui. Cette mission, lorsqu'elle est
exercée avec rigueur, intégrité et dans le respect des cadres légaux —
camerounais comme internationaux, est l'une des plus utiles que puisse
remplir un expert en cybersécurité.

\vspace{2ex}
\begin{center}
    {\color{CouleurAccent}\rule{0.6\linewidth}{1pt}}\\[0.8ex]
    {\small\itshape\color{CouleurPrimaire}
        Rédigé avec la conviction que la maîtrise technique ne prend de sens\\
        que lorsqu'elle est mise au service d'une éthique professionnelle irréprochable.}\\[0.5ex]
    {\small\bfseries\color{CouleurPrimaire}\auteurcours}\\[0.3ex]
    {\footnotesize\color{CouleurBordInfo}\itshape
        \institutionprimaire{} --- \anneescolaire}\\[0.8ex]
    {\color{CouleurAccent}\rule{0.6\linewidth}{1pt}}
\end{center}
